% ============================================================
%  BÁO CÁO ĐỒ ÁN CUỐI KHÓA — MÔN THỊ GIÁC MÁY TÍNH
%  ĐỀ TÀI: HỆ THỐNG PHÂN TÍCH TRẬN ĐẤU PICKLEBALL
%  SỬ DỤNG THỊ GIÁC MÁY TÍNH VÀ HỌC SÂU
% ============================================================
\documentclass[12pt,a4paper]{article}

% === PACKAGES ===
\usepackage[utf8]{inputenc}
\usepackage[vietnamese]{babel}
\usepackage{amsmath,amssymb}
\usepackage{graphicx}
\usepackage{booktabs}
\usepackage{tabularx}
\usepackage{float}
\usepackage{hyperref}
\usepackage{listings}
\usepackage{xcolor}
\usepackage{geometry}
\usepackage{caption}
\usepackage{subcaption}
\usepackage{algorithm}
\usepackage{algpseudocode}
\usepackage{tikz}
\usetikzlibrary{shapes,arrows,positioning,fit,calc}
\usepackage{enumitem}
\usepackage{fancyhdr}
\usepackage{setspace}
\usepackage{indentfirst}

\geometry{margin=2.5cm}
\onehalfspacing

% === CODE LISTING STYLE ===
\lstset{
    basicstyle=\small\ttfamily,
    keywordstyle=\color{blue}\bfseries,
    commentstyle=\color{gray},
    stringstyle=\color{red},
    breaklines=true,
    frame=single,
    numbers=left,
    numberstyle=\tiny\color{gray},
    language=Python,
    showstringspaces=false,
    captionpos=b,
}

% === HEADER/FOOTER ===
\pagestyle{fancy}
\fancyhead[L]{\small Đồ án cuối khóa --- Thị giác máy tính}
\fancyhead[R]{\small Trang~\thepage}
\fancyfoot[C]{}
\renewcommand{\headrulewidth}{0.4pt}

% ================================================================
% TRANG BÌA
% ================================================================
\begin{document}
\begin{titlepage}
\centering
\vspace*{1cm}

% Logo trường (uncomment và thay đường dẫn nếu có)
% \includegraphics[width=0.25\textwidth]{logo.png}\\[1cm]

{\large\textbf{TÊN TRƯỜNG ĐẠI HỌC}}\\[0.3cm]
{\large\textbf{KHOA CÔNG NGHỆ THÔNG TIN}}\\[2cm]

{\Huge\textbf{ĐỒ ÁN CUỐI KHÓA}}\\[0.5cm]
{\Large\textbf{MÔN: THỊ GIÁC MÁY TÍNH}}\\[2cm]

\rule{\textwidth}{1.5pt}\\[0.5cm]
{\LARGE\textbf{HỆ THỐNG PHÂN TÍCH TRẬN ĐẤU PICKLEBALL\\SỬ DỤNG THỊ GIÁC MÁY TÍNH VÀ HỌC SÂU}}\\[0.3cm]
{\large\textit{Pickleball Match Analysis System Using\\Computer Vision and Deep Learning}}\\[0.3cm]
\rule{\textwidth}{1.5pt}\\[2.5cm]

\begin{minipage}{0.45\textwidth}
\begin{flushleft}
\textbf{Giảng viên hướng dẫn:}\\
% Thay tên GVHD
TS. Nguyễn Văn A
\end{flushleft}
\end{minipage}
\hfill
\begin{minipage}{0.45\textwidth}
\begin{flushright}
\textbf{Sinh viên thực hiện:}\\
% Thay tên sinh viên
Nguyễn Văn B\\
MSSV: 20XXXXXX
\end{flushright}
\end{minipage}

\vfill
{\large Thành phố Hồ Chí Minh, \the\year}
\end{titlepage}

% ================================================================
% LỜI CẢM ƠN (tùy chọn -- xóa nếu không cần)
% ================================================================
\newpage
\section*{Lời cảm ơn}
\addcontentsline{toc}{section}{Lời cảm ơn}

% Thay nội dung lời cảm ơn
Em xin gửi lời cảm ơn chân thành đến giảng viên hướng dẫn đã tận tình hỗ trợ và định hướng trong suốt quá trình thực hiện đồ án. Em cũng xin cảm ơn các bạn trong lớp đã đóng góp ý kiến quý báu.

% ================================================================
% MỤC LỤC
% ================================================================
\newpage
\tableofcontents
\newpage
\listoffigures
\listoftables
\newpage

% ================================================================
% TÓM TẮT
% ================================================================
\begin{abstract}
Trong những năm gần đây, pickleball đã trở thành môn thể thao có tốc độ phát triển nhanh nhất tại Bắc Mỹ. Tuy nhiên, so với các môn thể thao khác như tennis hay bóng đá, pickleball vẫn còn thiếu vắng các hệ thống phân tích trận đấu tự động dựa trên thị giác máy tính. Báo cáo này trình bày việc thiết kế, triển khai và đánh giá một hệ thống phân tích trận đấu pickleball hoàn chỉnh từ video (end-to-end pipeline), kết hợp các kỹ thuật thị giác máy tính truyền thống và học sâu.

Hệ thống được đề xuất bao gồm ba thành phần chính: (i)~mô hình phát hiện 12 điểm đặc trưng sân dựa trên YOLO-Pose, phục vụ phép chiếu homography cục bộ theo vùng (zone-based local homography); (ii)~cơ chế phát hiện bóng cascade 4 tầng, kết hợp hai mô hình YOLO, pipeline thị giác máy tính truyền thống (phân đoạn màu HSV, phát hiện chuyển động, phân tích contour), và nội suy quỹ đạo đa thức; (iii)~mô hình phát hiện vận động viên với bộ lọc polygon sân và thuật toán phân loại đội tự động dựa trên vị trí lưới.

Ngoài ra, hệ thống tích hợp bộ lọc Kalman kết hợp Optical Flow để ổn định quỹ đạo bóng, thuật toán phát hiện bóng chạm sân (bounce detection) với bộ lọc phân biệt bounce và hit dựa trên phân tích vị trí tương đối bóng--vận động viên, cùng phân loại IN/OUT. Kết quả trực quan hóa được hiển thị trên hệ thống ba bản đồ sân 2D (triple minimap).

Kết quả thực nghiệm cho thấy mô hình phát hiện vận động viên đạt mAP@50 = 98,2\% (class Person) và 79,7\% (class Pickleball). Cơ chế cascade ball detection giúp tăng đáng kể tỷ lệ phát hiện bóng so với phương pháp đơn mô hình. Hệ thống được phát triển qua 7 phiên bản cải tiến liên tục, mỗi phiên bản giải quyết một nhóm vấn đề kỹ thuật cụ thể.

\vspace{0.3cm}
\noindent\textbf{Từ khóa:} Thị giác máy tính, Học sâu, Pickleball, YOLO, Phát hiện keypoint, Homography, Bộ lọc Kalman, Optical Flow, Phát hiện bóng chạm sân.
\end{abstract}

% ================================================================
\section{Giới thiệu}
\label{sec:intro}

\subsection{Đặt vấn đề}

Pickleball là môn thể thao vợt sử dụng bóng nhựa có lỗ, được chơi trên sân có kích thước $13{,}41 \times 6{,}10$~m (44~ft $\times$ 20~ft), nhỏ hơn đáng kể so với sân tennis tiêu chuẩn. Theo báo cáo của Sports \& Fitness Industry Association~\cite{sfia2023}, số lượng người chơi pickleball tại Hoa Kỳ đã tăng 158,6\% trong giai đoạn 2020--2023, đưa môn thể thao này trở thành hiện tượng phát triển nhanh nhất ngành thể thao Bắc Mỹ.

Sự bùng nổ này kéo theo nhu cầu lớn về công cụ phân tích trận đấu tự động, phục vụ ba mục đích chính: (i)~hỗ trợ huấn luyện viên đánh giá chiến thuật và kỹ năng vận động viên thông qua dữ liệu định lượng; (ii)~cung cấp thông tin bổ trợ cho trọng tài trong các tình huống tranh cãi (đặc biệt là phân loại IN/OUT); và (iii)~nâng cao trải nghiệm phát sóng với đồ họa phân tích thời gian thực.

Tuy nhiên, khác với tennis --- nơi hệ thống Hawk-Eye~\cite{hawkeye2024} sử dụng đa camera và tái tạo 3D đã được triển khai rộng rãi tại các giải Grand Slam --- pickleball hiện vẫn chưa có giải pháp phân tích tự động tương đương. Nguyên nhân chính xuất phát từ ba thách thức kỹ thuật đặc thù:

\begin{enumerate}[noitemsep]
    \item \textbf{Bóng nhỏ và tốc độ cao}: Bóng pickleball có đường kính khoảng 7,4~cm, chỉ chiếm 10--30 pixel trên khung hình 1080p. Tốc độ smash có thể vượt 100~km/h, gây hiện tượng motion blur nghiêm trọng.
    \item \textbf{Biến dạng phối cảnh lớn}: Sân nhỏ kết hợp với góc camera xiên tạo ra hiện tượng perspective warp mạnh, khiến phần sân xa camera bị nén đáng kể so với phần gần.
    \item \textbf{Nhiều vận động viên trên sân hẹp}: Với 4 vận động viên (doubles) trên sân nhỏ, các trường hợp che khuất (occlusion) xảy ra thường xuyên, gây khó khăn cho việc phát hiện và theo dõi.
\end{enumerate}

\subsection{Mục tiêu nghiên cứu}

Đồ án này đặt ra mục tiêu xây dựng một hệ thống phân tích trận đấu pickleball tự động từ video (end-to-end pipeline), bao gồm năm nhiệm vụ cụ thể:

\begin{enumerate}
    \item \textbf{Phát hiện sân}: Xác định vị trí 12 điểm đặc trưng (keypoints) trên sân pickleball và tính toán ma trận biến đổi homography để chiếu tọa độ camera sang tọa độ sân 2D.
    \item \textbf{Theo dõi bóng}: Phát hiện và theo dõi liên tục vị trí bóng pickleball qua các khung hình, sử dụng cơ chế cascade đa phương pháp để đạt tỷ lệ phát hiện cao.
    \item \textbf{Phát hiện vận động viên}: Xác định vị trí vận động viên trên sân, loại bỏ khán giả và trọng tài, đồng thời phân loại tự động vào hai đội dựa trên vị trí.
    \item \textbf{Phân tích bounce}: Phát hiện thời điểm bóng chạm sân, phân biệt với pha đánh bóng (hit), và phân loại IN/OUT.
    \item \textbf{Trực quan hóa}: Hiển thị kết quả phân tích trên hệ thống ba bản đồ sân 2D (triple minimap) tích hợp trực tiếp vào video đầu ra.
\end{enumerate}

\subsection{Phạm vi và giới hạn}

Hệ thống được thiết kế để hoạt động với video đơn camera (monocular), độ phân giải Full HD (1920$\times$1080), góc quay sideline hoặc overhead. Các giới hạn bao gồm: (i)~không hỗ trợ xử lý thời gian thực (offline processing); (ii)~chưa tích hợp player re-identification liên tục; và (iii)~độ chính xác phụ thuộc vào chất lượng phát hiện keypoints sân.

\subsection{Cấu trúc báo cáo}

Báo cáo được tổ chức như sau: Phần~\ref{sec:related} khảo sát các nghiên cứu liên quan. Phần~\ref{sec:methodology} trình bày phương pháp đề xuất bao gồm kiến trúc tổng quan và các thuật toán. Phần~\ref{sec:cv_techniques} mô tả chi tiết các kỹ thuật thị giác máy tính được sử dụng. Phần~\ref{sec:implementation} trình bày chi tiết triển khai. Phần~\ref{sec:experiments} phân tích kết quả thực nghiệm. Phần~\ref{sec:challenges} thảo luận khó khăn gặp phải và giải pháp. Phần~\ref{sec:conclusion} kết luận và đề xuất hướng phát triển.

% ================================================================
\section{Tổng quan nghiên cứu liên quan}
\label{sec:related}

\subsection{Phát hiện đối tượng dựa trên học sâu}

Bài toán phát hiện đối tượng (object detection) trong ảnh đã có bước tiến vượt bậc nhờ các kiến trúc mạng nơ-ron tích chập (CNN). Trong đó, họ kiến trúc YOLO (You Only Look Once)~\cite{yolo2016} đã trở thành tiêu chuẩn thực tiễn cho các bài toán yêu cầu phát hiện thời gian thực nhờ khả năng xử lý toàn bộ khung hình trong một lần truyền xuôi (single forward pass) duy nhất.

Từ YOLOv1 đến phiên bản mới nhất YOLO11~\cite{ultralytics2024}, kiến trúc đã liên tục được cải tiến về backbone (CSPDarknet, EfficientNet), neck (FPN, PANet), và head (anchor-free, decoupled head). YOLO11 medium (\texttt{yolo11m}) với 20 triệu tham số và 67,7 GFLOPs đạt mAP@50 = 51,5\% trên tập COCO~val, cân bằng tốt giữa tốc độ và độ chính xác, phù hợp cho các ứng dụng phân tích video thể thao.

\subsection{Phát hiện điểm đặc trưng sân thể thao}

Phương pháp phát hiện điểm đặc trưng sân (court keypoint detection) đã được nghiên cứu rộng rãi trong bối cảnh tennis. Faulkner và Dick~\cite{faulkner2017} đề xuất sử dụng mạng CNN để phát hiện trực tiếp tọa độ các đường kẻ sân tennis từ ảnh. Gần đây hơn, Maji và cộng sự~\cite{maji2022} mở rộng kiến trúc YOLO-Pose để phát hiện đồng thời bounding box và tọa độ keypoint, cho phép xác định vị trí sân chỉ qua một lần inference.

Cách tiếp cận dựa trên keypoint có nhiều ưu điểm so với phương pháp phát hiện đường thẳng (line detection) truyền thống sử dụng biến đổi Hough: (i)~output là tọa độ rời rạc cho từng điểm, tránh lỗi ghép nối đường; (ii)~mỗi keypoint có chỉ số confidence riêng, cho phép đánh giá độ tin cậy; và (iii)~keypoints cho phép tính toán phép chiếu homography cục bộ theo vùng.

\subsection{Theo dõi bóng trong thể thao}

Theo dõi bóng nhỏ tốc độ cao là bài toán thách thức trong thị giác máy tính thể thao. TrackNet~\cite{tracknet2019} đề xuất kiến trúc encoder-decoder chuyên biệt với đầu vào là 3 khung hình liên tiếp, đạt kết quả ấn tượng trong theo dõi bóng badminton và tennis. Tuy nhiên, TrackNet yêu cầu huấn luyện chuyên dụng cho từng môn thể thao.

Các phương pháp kết hợp bộ lọc Kalman~\cite{kalman1960} với Optical Flow (dòng quang học)~\cite{lucas1981} đã chứng minh hiệu quả trong theo dõi đối tượng qua các khung hình bị thiếu phát hiện (missed detection). Bộ lọc Kalman cung cấp ước tính vị trí tối ưu dựa trên mô hình chuyển động, trong khi Optical Flow cho phép ước tính vận tốc đối tượng từ sự thay đổi cường độ pixel giữa hai khung hình liên tiếp.

\subsection{Phép biến đổi phối cảnh (Homography)}

Phép biến đổi homography~\cite{hartley2003} là công cụ toán học cơ bản để ánh xạ giữa hai mặt phẳng trong không gian chiếu. Trong phân tích thể thao, homography được sử dụng để ánh xạ tọa độ đối tượng từ mặt phẳng camera sang mặt phẳng sân (bird's-eye view). Tuy nhiên, phép biến đổi này có một giới hạn cố hữu: chỉ chính xác cho các điểm nằm trên mặt phẳng tham chiếu (mặt sân). Khi áp dụng cho đối tượng bay trên không (bóng đang bay), sai số chiếu tăng tỷ lệ thuận với độ cao đối tượng so với mặt sân.

\subsection{Phân tích trận đấu pickleball --- Tình hình hiện tại}

So với tennis và bóng đá, nghiên cứu về phân tích tự động trận đấu pickleball còn rất hạn chế. Các dự án mã nguồn mở hiện có chủ yếu tập trung vào từng thành phần đơn lẻ (chỉ theo dõi bóng hoặc chỉ phát hiện sân), chưa có pipeline tích hợp toàn diện. Đồ án này đóng góp vào lĩnh vực bằng việc xây dựng hệ thống end-to-end kết hợp cả ba thành phần (sân, bóng, vận động viên) cùng với phân tích bounce và trực quan hóa.

% ============================================================
\section{Phương pháp đề xuất}
\label{sec:methodology}

\subsection{Tổng quan kiến trúc}

\begin{figure}[H]
\centering
\begin{tikzpicture}[
    node distance=0.8cm,
    block/.style={rectangle, draw, fill=blue!10, text width=5cm, text centered, rounded corners, minimum height=0.8cm, font=\small},
    arrow/.style={->, >=stealth, thick},
]
    \node[block, fill=green!10] (input) {Video Input (1920$\times$1080)};
    \node[block, below=of input] (court) {Court Keypoint Detection \\ \texttt{court\_keypoint\_best.pt}};
    \node[block, below=of court] (ball) {Cascade Ball Detection \\ (4 tầng)};
    \node[block, below=of ball] (player) {Player Detection + Team \\ \texttt{player\_detection\_best.pt}};
    \node[block, below=of player, fill=orange!10] (proj) {Zone-Based Homography \\ Projection};
    \node[block, below=of proj, fill=red!10] (bounce) {Bounce Detection \\ + Hit Filter};
    \node[block, below=of bounce, fill=yellow!10] (viz) {Triple Minimap Visualization};

    \draw[arrow] (input) -- (court);
    \draw[arrow] (input.east) -- ++(1.5,0) |- (ball.east);
    \draw[arrow] (input.west) -- ++(-1.5,0) |- (player.west);
    \draw[arrow] (court) -- (ball);
    \draw[arrow] (ball) -- (player);
    \draw[arrow] (player) -- (proj);
    \draw[arrow] (court.east) -- ++(2.3,0) |- (proj.east);
    \draw[arrow] (proj) -- (bounce);
    \draw[arrow] (bounce) -- (viz);
\end{tikzpicture}
\caption{Kiến trúc tổng quan của pipeline phân tích.}
\label{fig:architecture}
\end{figure}

\subsection{Phát hiện sân bằng Keypoint}
\label{sec:court_detection}

Thay vì phương pháp phát hiện đường thẳng (line detection) truyền thống, hệ thống sử dụng mô hình YOLO-Pose để phát hiện \textbf{12 điểm đặc trưng} (keypoints) trên sân. Mỗi keypoint $k_i = (x_i, y_i, c_i)$ bao gồm tọa độ pixel $(x_i, y_i)$ và confidence $c_i \in [0, 1]$.

\subsubsection{Dataset}
\begin{table}[H]
\centering
\caption{Thông tin dataset training --- Court Keypoint Model}
\label{tab:dataset_court}
\begin{tabularx}{\textwidth}{l X}
\toprule
\textbf{Thuộc tính} & \textbf{Giá trị} \\
\midrule
Nguồn & Roboflow Universe --- ``pb-9bsin'' (PB Workspace) \\
URL & \url{https://universe.roboflow.com/pb-rrerm/pb-9bsin/dataset/5} \\
Số ảnh & 2,064 ảnh gốc (đã chia train/val/test, có augmentation) \\
Task & Keypoint Detection --- 12 keypoints trên sân \\
License & CC BY 4.0 \\
Lý do chọn & Dataset duy nhất trên Roboflow có annotation 12 keypoints sân pickleball, đã chia sẵn và có augmentation \\
\bottomrule
\end{tabularx}
\end{table}

\subsubsection{Cấu hình Training}
\begin{table}[H]
\centering
\caption{Cấu hình huấn luyện --- Court Keypoint Model}
\label{tab:training_court}
\begin{tabular}{ll|ll}
\toprule
\textbf{Tham số} & \textbf{Giá trị} & \textbf{Tham số} & \textbf{Giá trị} \\
\midrule
Model & YOLOv8m-pose & Epochs & 100 \\
GPU & NVIDIA A100-SXM4-40GB & Batch size & 16 \\
Optimizer & AdamW & Learning rate & 0.001 \\
Image size & 640$\times$640 & AMP & FP16 \\
Task & Pose (keypoint) & Keypoints & 12 \\
Mosaic & 1.0 & Patience & 20 (early stop) \\
\bottomrule
\end{tabular}
\end{table}

\subsubsection{Lý do chuyển từ Line Detection sang Keypoint Detection}
\begin{itemize}[noitemsep]
    \item Line detection yêu cầu nối các đoạn thẳng theo thứ tự đúng --- lỗi thứ tự tạo đường chéo cắt ngang sân
    \item Keypoint detection cho output riêng biệt cho từng điểm, tránh lỗi nối
    \item Keypoints cho phép tính homography cục bộ theo vùng (zone-based)
\end{itemize}

\subsubsection{Bố trí 12 keypoints}
12 keypoints được sắp xếp trên sân theo cấu trúc 4 hàng $\times$ 3 cột:

\begin{figure}[H]
\centering
\begin{tikzpicture}[scale=0.7]
    % Court outline
    \draw[thick] (0,0) rectangle (6,8.8);
    \draw[thick] (0,3) -- (6,3);
    \draw[thick] (0,5.8) -- (6,5.8);
    \draw[thick, dashed] (3,0) -- (3,3);
    \draw[thick, dashed] (3,5.8) -- (3,8.8);
    \draw[very thick, red] (0,4.4) -- (6,4.4);
    \node[right, red] at (6,4.4) {\small NET};

    % Keypoints with zigzag numbering
    \foreach \x/\y/\label in {0/8.8/0, 3/8.8/1, 6/8.8/2,
                               6/5.8/3, 3/5.8/4, 0/5.8/5,
                               0/3/6, 3/3/7, 6/3/8,
                               6/0/9, 3/0/10, 0/0/11} {
        \fill[blue] (\x,\y) circle (4pt);
        \node[above right, font=\scriptsize\bfseries] at (\x,\y) {\label};
    }

    % Zone labels
    \node at (1.5,7.3) {\small Zone A};
    \node at (4.5,7.3) {\small Zone B};
    \node at (1.5,4.9) {\small Zone C};
    \node at (4.5,4.9) {\small Zone D};
    \node at (1.5,1.5) {\small Zone E};
    \node at (4.5,1.5) {\small Zone F};

    % Arrows showing zigzag
    \draw[->, thick, gray] (-0.5,8.8) -- (-0.5,8.3) node[left, font=\tiny] {L$\to$R};
    \draw[<-, thick, gray] (-0.5,5.8) -- (-0.5,6.3) node[left, font=\tiny] {R$\to$L};
    \draw[->, thick, gray] (-0.5,3) -- (-0.5,2.5) node[left, font=\tiny] {L$\to$R};
    \draw[<-, thick, gray] (-0.5,0) -- (-0.5,0.5) node[left, font=\tiny] {R$\to$L};
\end{tikzpicture}
\caption{Bố trí 12 keypoints trên sân và hướng đánh số zigzag của mô hình.}
\label{fig:keypoints}
\end{figure}

\subsubsection{Phát hiện thứ tự Zigzag}
\label{sec:zigzag}
Trong quá trình thực nghiệm, phát hiện mô hình đánh số keypoints theo thứ tự \textbf{zigzag} (serpentine): hàng chẵn (0, 2) đánh số trái$\to$phải, hàng lẻ (1, 3) đánh số phải$\to$trái. Điều này khiến mapping tọa độ ban đầu (giả sử tất cả hàng đi trái$\to$phải) bị \textbf{gương ngược trái/phải}, dẫn đến vận động viên chiếu sai phía trên minimap.

Giải pháp: đảo tọa độ $x$ trong mảng tọa độ đích (COURT\_DST) cho hàng 1 và 3:

\begin{equation}
\text{COURT\_DST}[i] = \begin{cases}
    (x, y) & \text{nếu hàng chẵn (đánh số L}\to\text{R)} \\
    (W - x, y) & \text{nếu hàng lẻ (đánh số R}\to\text{L)}
\end{cases}
\label{eq:zigzag}
\end{equation}

trong đó $W = 400$ là chiều rộng sân chuẩn hóa.

\subsection{Chiếu phối cảnh theo vùng (Zone-Based Projection)}
\label{sec:zone_projection}

\subsubsection{Homography toàn cục vs cục bộ}
Một ma trận homography toàn cục $H_{global}$ ánh xạ toàn bộ sân, nhưng sai số tăng ở các vùng xa các keypoint tham chiếu. Để giảm sai số, hệ thống chia sân thành \textbf{6 vùng} (zones A--F), mỗi vùng có homography cục bộ riêng.

\subsubsection{Tính toán}
Cho mỗi zone $Z$ với 4 keypoints góc $\{k_{i_1}, k_{i_2}, k_{i_3}, k_{i_4}\}$:

\begin{equation}
H_Z = \texttt{cv2.getPerspectiveTransform}(S_Z, D_Z)
\end{equation}

trong đó $S_Z$ là tọa độ camera-space và $D_Z$ là tọa độ sân 2D tương ứng. Khi chiếu điểm $p$:

\begin{equation}
p_{2D} = \begin{cases}
    H_Z \cdot p & \text{nếu } p \in Z \\
    H_{global} \cdot p & \text{nếu } p \notin \text{bất kỳ zone nào}
\end{cases}
\end{equation}

\subsubsection{Kiểm tra độ tin cậy}
Khi camera zoom cận (close-up, replay), keypoints vẫn được phát hiện nhưng ở vị trí sai. Hàm \texttt{is\_reliable()} kiểm tra 3 tiêu chí:
\begin{enumerate}[noitemsep]
    \item Số keypoints có $c_i > 0.3$ phải $\geq 6$
    \item Confidence trung bình phải $\geq 0.4$
    \item Diện tích tứ giác 4 góc sân phải $\geq 5000$ px\textsuperscript{2}
\end{enumerate}

% ============================================================
\subsection{Cascade Ball Detection}
\label{sec:cascade_ball}

Bóng pickleball có đường kính $\sim$7cm, chỉ chiếm 10--30 pixel trên frame 1080p, và thường bị blur khi bay nhanh. Một mô hình duy nhất cho detection rate thấp.

\subsubsection{Dataset}
\begin{table}[H]
\centering
\caption{Thông tin dataset training --- Ball Detection Model}
\label{tab:dataset_ball}
\begin{tabularx}{\textwidth}{l X}
\toprule
\textbf{Thuộc tính} & \textbf{Giá trị} \\
\midrule
Nguồn & Roboflow Universe --- ``annotations\_pickleball'' (Pickleballbadminton Workspace) \\
URL & \url{https://universe.roboflow.com/pickleballbadminton/annotations_pickleball/dataset/8} \\
Số ảnh & 18,408 ảnh bóng pickleball \\
Classes & Pickleball (1 class --- chỉ bóng) \\
License & CC BY 4.0 \\
Lý do chọn & Dataset lớn nhất cho bóng pickleball (18K+ ảnh), đa dạng góc camera và điều kiện ánh sáng, giúp model generalize tốt \\
\bottomrule
\end{tabularx}
\end{table}

\subsubsection{Cấu hình Training}
\begin{table}[H]
\centering
\caption{Cấu hình huấn luyện --- Ball Detection Model}
\label{tab:training_ball}
\begin{tabular}{ll|ll}
\toprule
\textbf{Tham số} & \textbf{Giá trị} & \textbf{Tham số} & \textbf{Giá trị} \\
\midrule
Model & YOLOv8n (nano) & Epochs & 100 \\
GPU & NVIDIA T4 16GB & Batch size & 64 \\
Optimizer & AdamW & Learning rate & 0.01 \\
Image size & 640$\times$640 & AMP & FP16 \\
Classes & 1 (Pickleball) & Patience & 15 (early stop) \\
Mosaic & 1.0 & Mixup & 0.15 \\
\bottomrule
\end{tabular}
\end{table}

\subsubsection{Cơ chế Cascade 4 tầng}
Hệ thống sử dụng cơ chế \textbf{cascade 4 tầng}:

\begin{table}[H]
\centering
\caption{Cascade Ball Detection --- 4 tầng tuần tự}
\label{tab:cascade}
\begin{tabularx}{\textwidth}{c l X c}
\toprule
\textbf{Tầng} & \textbf{Phương pháp} & \textbf{Chi tiết} & \textbf{Conf.} \\
\midrule
1 & YOLO Ball Model & Mô hình chuyên biệt cho bóng pickleball (\texttt{ball\_tracker\_best.pt}) & 0.15 \\
2 & YOLO Player Model & Class ``Pickleball'' từ mô hình player (\texttt{player\_detection\_best.pt}) & 0.20 \\
3 & Classical CV & HSV color filter + motion mask + contour circularity & --- \\
4 & Trajectory Interp. & Polynomial degree-2 fit, tối đa 5 frame dự đoán & --- \\
\bottomrule
\end{tabularx}
\end{table}

\subsubsection{Tầng 3: Classical Computer Vision}
Khi cả hai mô hình YOLO không phát hiện được, hệ thống sử dụng pipeline CV truyền thống:
\begin{enumerate}[noitemsep]
    \item \textbf{Color segmentation}: Lọc vùng HSV$(20\text{--}55, 80\text{--}255, 150\text{--}255)$ --- màu vàng/xanh lá của bóng
    \item \textbf{Motion mask}: Hiệu khung hình liên tiếp $|I_t - I_{t-1}|$ --- loại vật thể tĩnh
    \item \textbf{Morphology}: Open + Close bằng kernel ellipse $3\times3$
    \item \textbf{Contour analysis}: Lọc theo diện tích ($4 \leq A \leq 800$ px\textsuperscript{2}) và circularity $\geq 0.3$
\end{enumerate}

\subsubsection{Tầng 4: Trajectory Interpolation}
Khi tất cả phương pháp phát hiện thất bại, hệ thống sử dụng nội suy đa thức:

\begin{equation}
\hat{x}(t) = \sum_{k=0}^{d} a_k t^k, \quad \hat{y}(t) = \sum_{k=0}^{d} b_k t^k
\label{eq:poly}
\end{equation}

với $d = 2$ (bậc 2) và $t$ là chỉ số frame. Các hệ số $a_k, b_k$ được fit từ 15 frame gần nhất bằng phương pháp bình phương tối thiểu. Dự đoán tối đa 5 frame liên tiếp trước khi dừng.

\subsection{Kalman Filter kết hợp Optical Flow}
\label{sec:kalman_of}

\subsubsection{Mô hình trạng thái}
Bộ lọc Kalman với vector trạng thái $\mathbf{x} = [x, y, v_x, v_y]^T$:

\begin{equation}
\mathbf{x}_{t+1} = \mathbf{F} \cdot \mathbf{x}_t + \mathbf{w}_t, \quad \mathbf{z}_t = \mathbf{H} \cdot \mathbf{x}_t + \mathbf{v}_t
\end{equation}

với ma trận chuyển trạng thái và quan sát:

\begin{equation}
\mathbf{F} = \begin{bmatrix} 1 & 0 & \Delta t & 0 \\ 0 & 1 & 0 & \Delta t \\ 0 & 0 & 1 & 0 \\ 0 & 0 & 0 & 1 \end{bmatrix}, \quad
\mathbf{H} = \begin{bmatrix} 1 & 0 & 0 & 0 \\ 0 & 1 & 0 & 0 \end{bmatrix}
\end{equation}

\subsubsection{Tích hợp Optical Flow}
Khi ball detector không phát hiện được, vận tốc từ Lucas-Kanade Optical Flow \cite{lucas1981} được blend vào trạng thái Kalman:

\begin{equation}
v_{x}^{blend} = (1 - \alpha) \cdot v_{x}^{KF} + \alpha \cdot v_{x}^{OF}, \quad \alpha = 0.3
\end{equation}

Phương pháp này duy trì tracking mượt mà qua các frame không có detection.

\subsubsection{Ma trận nhiễu vật lý}
Ma trận nhiễu quá trình $\mathbf{Q}$ sử dụng mô hình vật lý thay vì giá trị cố định:

\begin{equation}
\mathbf{Q} = \sigma_q^2 \begin{bmatrix}
    \frac{\Delta t^4}{4} & 0 & \frac{\Delta t^3}{2} & 0 \\
    0 & \frac{\Delta t^4}{4} & 0 & \frac{\Delta t^3}{2} \\
    \frac{\Delta t^3}{2} & 0 & \Delta t^2 & 0 \\
    0 & \frac{\Delta t^3}{2} & 0 & \Delta t^2
\end{bmatrix}
\end{equation}

% ============================================================
\subsection{Phát hiện vận động viên}
\label{sec:player_detection}

\subsubsection{Vì sao cần mô hình chuyên biệt?}
Mặc dù YOLO pretrained trên COCO đã có class ``person'', mô hình chuyên biệt cần thiết vì:
\begin{enumerate}[noitemsep]
    \item \textbf{Không phân biệt vai trò}: COCO ``person'' phát hiện tất cả người trong khung hình --- bao gồm khán giả, trọng tài, cameraman --- gây nhiễu
    \item \textbf{Context-specific accuracy}: Mô hình train trên ảnh sân pickleball cho accuracy cao hơn trong ngữ cảnh cụ thể
    \item \textbf{Dual-purpose}: Mô hình custom detect cả 2 classes (Person + Pickleball), phục vụ làm tầng 2 trong cascade ball detection
\end{enumerate}

\subsubsection{Dataset}
\begin{table}[H]
\centering
\caption{Thông tin dataset training}
\label{tab:dataset}
\begin{tabularx}{\textwidth}{l X}
\toprule
\textbf{Thuộc tính} & \textbf{Giá trị} \\
\midrule
Nguồn & Roboflow Universe --- ``Pickleball with Players'' (Hughs Workspace) \\
URL & \url{https://universe.roboflow.com/hughs-workspace-plw3g/pickleball-with-players-topw1} \\
Số ảnh & 4,119 (train: 6,710 / val: 564 / test: có) \\
Classes & Person, Pickleball (2 classes) \\
License & CC BY 4.0 \\
Lý do chọn & Lớn nhất trong các dataset pickleball trên Roboflow, có cả player và ball, đa dạng góc camera \\
\bottomrule
\end{tabularx}
\end{table}

\subsubsection{Cấu hình Training}
\begin{table}[H]
\centering
\caption{Hyperparameters training player detection model}
\label{tab:training}
\begin{tabular}{ll|ll}
\toprule
\textbf{Param} & \textbf{Value} & \textbf{Param} & \textbf{Value} \\
\midrule
Model & YOLO11m (20M params) & Epochs & 80 \\
GPU & NVIDIA A100 40GB & Batch size & 32 \\
Optimizer & AdamW & Learning rate & 0.001 \\
Image size & 640$\times$640 & AMP & FP16 \\
Mosaic & 1.0 & Mixup & 0.1 \\
Copy-paste & 0.1 & Patience & 15 (early stop) \\
\bottomrule
\end{tabular}
\end{table}

\subsubsection{Lọc người ngoài sân (Court Polygon Filter)}
Sử dụng 4 keypoints góc sân ($k_0, k_2, k_9, k_{11}$) tạo polygon mở rộng 15\%:

\begin{equation}
p'_i = p_i + 0.15 \cdot (p_i - \bar{p})
\end{equation}

trong đó $\bar{p}$ là tâm polygon. Hàm \texttt{cv2.pointPolygonTest()} kiểm tra điểm chân (foot point) của mỗi player có nằm trong polygon hay không.

\subsubsection{Phân loại đội (Team Assignment)}
Ban đầu sử dụng phương pháp ép chia 2+2 (2 gần + 2 xa camera) --- gây lỗi khi chỉ phát hiện được 3 vận động viên (1 người xa bị gán nhầm vào đội gần).

Giải pháp: sử dụng \textbf{vị trí lưới (Net Y)} từ keypoints làm ranh giới:

\begin{equation}
\text{Net}_Y = \frac{1}{|\mathcal{K}_{net}|} \sum_{i \in \mathcal{K}_{net}} y_i, \quad \mathcal{K}_{net} = \{3,4,5,6,7,8\}
\end{equation}

\begin{equation}
\text{Team}(p) = \begin{cases}
    A \text{ (Near)} & \text{nếu } \text{foot}_Y(p) \geq \text{Net}_Y \\
    B \text{ (Far)}  & \text{nếu } \text{foot}_Y(p) < \text{Net}_Y
\end{cases}
\end{equation}

% ============================================================
\subsection{Phát hiện bóng chạm sân (Bounce Detection)}
\label{sec:bounce_detection}

\subsubsection{Phương pháp phát hiện}
Bounce được xác định khi tọa độ $Y$ (camera-space) của bóng đạt \textbf{cực đại cục bộ} trong cửa sổ trượt $w = 7$ frame:

\begin{equation}
\text{isBounce}(t) = \begin{cases}
    \text{True} & \text{nếu } y_t > \max(y_{t-w/2:t}) - 3 \text{ và } y_t > \max(y_{t:t+w/2}) - 3 \\
    \text{False} & \text{ngược lại}
\end{cases}
\end{equation}

với khoảng cách tối thiểu giữa 2 bounce là $\text{min\_gap} = 8$ frame.

\subsubsection{Bộ lọc phân biệt Bounce và Hit}
\label{sec:bounce_hit}
Vấn đề: cả bounce (chạm sân) và hit (vận động viên đánh bóng) đều làm bóng đổi hướng $Y$. Bộ lọc phân biệt dựa trên vị trí tương đối:

\begin{equation}
\text{isHit}(b, P) = \exists p \in P: \begin{cases}
    |b_x - \text{center}_x(p)| \leq w_p/2 + 30 & \text{(gần ngang)} \\
    b_y < \text{foot}_y(p) - 0.35 \cdot h_p & \text{(cao ngang thân)}
\end{cases}
\end{equation}

Khi bóng ở sát chân ($b_y \geq \text{foot}_y - 0.1 \cdot h_p$) vẫn được tính là bounce. Vùng $0.1h_p$ đến $0.35h_p$ trên foot là \textit{ambiguous zone} --- bỏ qua để tránh false positive.

\subsubsection{Phân loại IN/OUT}
Tại thời điểm bounce, vị trí 2D được kiểm tra:

\begin{equation}
\text{IN/OUT} = \begin{cases}
    \text{IN} & \text{nếu } -5 \leq x_{2D} \leq W+5 \text{ và } -5 \leq y_{2D} \leq H+5 \\
    \text{OUT} & \text{ngược lại}
\end{cases}
\end{equation}

với $W = 400, H = 880$ là kích thước sân chuẩn hóa và margin $= 5$ pixel cho sai số projection.

% ============================================================
\section{Nền tảng kỹ thuật Thị giác Máy tính}
\label{sec:cv_techniques}

Thị giác máy tính (Computer Vision) là nền tảng cốt lõi xuyên suốt toàn bộ hệ thống. Mọi thành phần trong pipeline --- từ tiền xử lý ảnh, phát hiện đối tượng, theo dõi chuyển động, đến chiếu phối cảnh và trực quan hóa --- đều được xây dựng trên các nguyên lý và thuật toán thị giác máy tính. Trong đó, học sâu (deep learning) đóng vai trò là \textit{một trong những công cụ} của thị giác máy tính hiện đại, bên cạnh các kỹ thuật truyền thống vẫn giữ vai trò thiết yếu mà không thể thay thế bằng mạng nơ-ron.

Phần này trình bày chi tiết các kỹ thuật thị giác máy tính được sử dụng, bao gồm: xử lý không gian màu, phân đoạn ảnh, phát hiện chuyển động, phép toán hình thái học, phân tích contour, phép biến đổi phối cảnh (homography), dòng quang học (optical flow), kiểm tra điểm trong đa giác (point-in-polygon), và kỹ thuật render bản đồ 2D. Mỗi kỹ thuật đóng vai trò không thể thiếu trong ít nhất một giai đoạn của pipeline.

\subsection{Tiền xử lý ảnh (Image Preprocessing)}

Mỗi frame video đầu vào ($1920 \times 1080$, BGR) trải qua các bước tiền xử lý trước khi đưa vào các module phân tích:

\begin{enumerate}[noitemsep]
    \item \textbf{Resize cho YOLO}: Ảnh được scale về $640 \times 640$ với letterboxing (thêm viền đen) để giữ tỷ lệ khung hình, tránh biến dạng đối tượng
    \item \textbf{Chuyển đổi color space}: BGR $\to$ RGB (cho YOLO), BGR $\to$ HSV (cho classical detection), BGR $\to$ Grayscale (cho optical flow và motion detection)
    \item \textbf{Half Precision (FP16)}: Pixel values được convert từ FP32 sang FP16 trước khi inference, giảm 50\% VRAM và tăng tốc trên GPU có Tensor Cores (A100, T4)
\end{enumerate}

\begin{equation}
I_{FP16}(x,y) = \text{float16}\left(\frac{I_{uint8}(x,y)}{255.0}\right)
\end{equation}

\subsection{Không gian màu HSV và Color Segmentation}
\label{sec:hsv}

Bóng pickleball có màu vàng/xanh lá đặc trưng. Thay vì lọc trong không gian BGR (phụ thuộc nhiều vào ánh sáng), hệ thống sử dụng \textbf{không gian màu HSV} (Hue-Saturation-Value) vì:

\begin{itemize}[noitemsep]
    \item \textbf{Hue (H)}: Mã hóa màu sắc thuần túy, ít bị ảnh hưởng bởi bóng đổ và ánh sáng
    \item \textbf{Saturation (S)}: Phân biệt màu rực (bóng) vs màu nhạt (sân)
    \item \textbf{Value (V)}: Loại vùng quá tối (bóng đổ) hoặc quá sáng (phản chiếu)
\end{itemize}

Chuyển đổi từ BGR sang HSV:
\begin{equation}
H = \begin{cases}
    60 \cdot \frac{G-B}{\Delta} & \text{nếu } V = R \\
    60 \cdot \left(\frac{B-R}{\Delta} + 2\right) & \text{nếu } V = G \\
    60 \cdot \left(\frac{R-G}{\Delta} + 4\right) & \text{nếu } V = B
\end{cases}
\end{equation}

trong đó $V = \max(R,G,B)$ và $\Delta = V - \min(R,G,B)$.

\textbf{Ngưỡng lọc bóng pickleball}: HSV $\in [20, 80, 150] \to [55, 255, 255]$

\begin{equation}
M_{color}(x,y) = \begin{cases}
    255 & \text{nếu } H_{low} \leq H(x,y) \leq H_{high} \text{ và } S_{low} \leq S \leq S_{high} \text{ và } V_{low} \leq V \leq V_{high} \\
    0 & \text{ngược lại}
\end{cases}
\end{equation}

Kết quả là binary mask $M_{color}$ chứa vùng có màu vàng/xanh lá --- ứng cử viên là bóng.

\subsection{Phát hiện chuyển động (Motion Detection)}
\label{sec:motion}

Mask màu sắc có thể chứa nhiều false positive (logo vàng, quần áo, biển quảng cáo). Để loại bỏ \textbf{vật thể tĩnh}, hệ thống sử dụng kỹ thuật \textbf{hiệu khung hình liên tiếp} (frame differencing):

\begin{equation}
D(x,y) = |I_t^{gray}(x,y) - I_{t-1}^{gray}(x,y)|
\end{equation}

\begin{equation}
M_{motion}(x,y) = \begin{cases}
    255 & \text{nếu } D(x,y) > \tau_{motion} \\
    0 & \text{ngược lại}
\end{cases}
\end{equation}

với $\tau_{motion} = 25$ (ngưỡng chuyển động). Mask chuyển động được \textbf{giãn nở} (dilate) 2 lần để kết nối các vùng phân mảnh:

\begin{equation}
M_{motion}' = M_{motion} \oplus K_{3\times3} \oplus K_{3\times3}
\end{equation}

Kết hợp hai mask bằng phép AND:
\begin{equation}
M_{final} = M_{color} \cap M_{motion}'
\end{equation}

Kết quả: chỉ giữ vùng \textit{vừa có màu vàng/xanh} VÀ \textit{vừa đang di chuyển} --- loại bỏ hiệu quả logo, biển quảng cáo tĩnh.

\subsection{Phép toán hình thái học (Morphological Operations)}
\label{sec:morphology}

Sau khi kết hợp mask, vẫn còn nhiễu dạng điểm nhỏ và lỗ hổng. Hệ thống áp dụng hai phép toán hình thái học sử dụng \textbf{structuring element hình elip} $K_{3\times3}$:

\subsubsection{Opening (Erosion $\to$ Dilation)}
Loại bỏ nhiễu nhỏ (noise dots):
\begin{equation}
M_{open} = (M_{final} \ominus K) \oplus K
\end{equation}

\textit{Erosion} co nhỏ tất cả vùng trắng --- vùng quá nhỏ (nhiễu) biến mất. \textit{Dilation} sau đó phục hồi kích thước vùng lớn.

\subsubsection{Closing (Dilation $\to$ Erosion)}
Lấp lỗ hổng bên trong vùng phát hiện:
\begin{equation}
M_{clean} = (M_{open} \oplus K) \ominus K
\end{equation}

\textit{Dilation} trước tiên lấp các lỗ nhỏ bên trong vùng bóng. \textit{Erosion} sau đó khôi phục viền chính xác.

\begin{figure}[H]
\centering
\begin{tikzpicture}[
    box/.style={rectangle, draw, minimum width=2cm, minimum height=0.8cm, text centered, font=\small},
    arrow/.style={->, >=stealth, thick},
]
    \node[box, fill=yellow!20] (color) {HSV Mask};
    \node[box, fill=green!20, right=1cm of color] (motion) {Motion Mask};
    \node[box, below=0.5cm of $(color)!0.5!(motion)$] (and) {AND};
    \node[box, below=0.5cm of and, fill=orange!20] (open) {Opening};
    \node[box, below=0.5cm of open, fill=orange!20] (close) {Closing};
    \node[box, below=0.5cm of close, fill=red!20] (contour) {Contour Analysis};
    
    \draw[arrow] (color) -- (and);
    \draw[arrow] (motion) -- (and);
    \draw[arrow] (and) -- (open);
    \draw[arrow] (open) -- (close);
    \draw[arrow] (close) -- (contour);
\end{tikzpicture}
\caption{Pipeline xử lý ảnh trong Classical Ball Detection (tầng 3).}
\label{fig:cv_pipeline}
\end{figure}

\subsection{Phân tích contour (Contour Analysis)}
\label{sec:contour}

Trên mask đã làm sạch, hệ thống tìm các \textbf{contour} (đường viền) bằng thuật toán Suzuki-Abe \cite{suzuki1985} thông qua \texttt{cv2.findContours()}. Mỗi contour được đánh giá theo 3 tiêu chí:

\subsubsection{Diện tích (Area)}
\begin{equation}
A = \sum_{i=0}^{n-1} \frac{1}{2} |x_i(y_{i+1} - y_{i-1})|
\label{eq:area}
\end{equation}

Ngưỡng: $4 \leq A \leq 800$ px$^2$. Quá nhỏ = nhiễu. Quá lớn = không phải bóng (có thể là tay VĐV, vật khác).

\subsubsection{Chu vi (Perimeter)}
Tính bằng tổng khoảng cách giữa các điểm liên tiếp trên contour:
\begin{equation}
P = \sum_{i=0}^{n-1} \sqrt{(x_{i+1} - x_i)^2 + (y_{i+1} - y_i)^2}
\end{equation}

\subsubsection{Circularity (Độ tròn)}
Đánh giá mức độ tròn của contour --- bóng pickleball là hình cầu nên contour phải gần tròn:
\begin{equation}
\text{Circularity} = \frac{4\pi A}{P^2}
\label{eq:circularity}
\end{equation}

Giá trị = 1.0 cho hình tròn hoàn hảo, < 0.3 cho hình dài/méo. Ngưỡng: Circularity $\geq 0.3$.

\subsubsection{Tâm contour (Centroid)}
Vị trí bóng = tâm contour, tính bằng \textbf{image moments}:
\begin{equation}
c_x = \frac{m_{10}}{m_{00}}, \quad c_y = \frac{m_{01}}{m_{00}}
\end{equation}

trong đó $m_{pq} = \sum_{x} \sum_{y} x^p y^q \cdot I(x,y)$ là spatial moment bậc $(p, q)$.

\subsection{Phép biến đổi phối cảnh (Perspective Transform)}
\label{sec:perspective}

\subsubsection{Cơ sở toán học}
Phép biến đổi homography ánh xạ điểm trên mặt phẳng camera sang mặt phẳng sân thông qua ma trận $3 \times 3$:

\begin{equation}
\begin{bmatrix} x' \\ y' \\ w \end{bmatrix} = 
\mathbf{H} \cdot \begin{bmatrix} x \\ y \\ 1 \end{bmatrix} =
\begin{bmatrix} h_{11} & h_{12} & h_{13} \\ h_{21} & h_{22} & h_{23} \\ h_{31} & h_{32} & h_{33} \end{bmatrix}
\begin{bmatrix} x \\ y \\ 1 \end{bmatrix}
\end{equation}

Tọa độ 2D sân cuối cùng:
\begin{equation}
x_{2D} = \frac{x'}{w} = \frac{h_{11}x + h_{12}y + h_{13}}{h_{31}x + h_{32}y + h_{33}}, \quad
y_{2D} = \frac{y'}{w} = \frac{h_{21}x + h_{22}y + h_{23}}{h_{31}x + h_{32}y + h_{33}}
\end{equation}

\subsubsection{Global vs Local Homography}
\textbf{Global}: Sử dụng tất cả 12 keypoints, tính $\mathbf{H}$ bằng RANSAC qua \texttt{cv2.findHomography()}, yêu cầu $\geq 4$ cặp điểm tương ứng. Sai số trung bình thấp nhưng sai số \textit{cục bộ} có thể cao ở vùng xa các keypoint.

\textbf{Local (Zone-Based)}: Sử dụng chính xác 4 keypoints góc mỗi zone, tính $\mathbf{H}_Z$ bằng \texttt{cv2.getPerspectiveTransform()} --- giải hệ 8 ẩn (8 phương trình từ 4 cặp điểm), cho kết quả exact (không cần RANSAC). Sai số cục bộ thấp hơn đáng kể vì 4 điểm tham chiếu bao quanh vùng chiếu.

\subsubsection{Xác định zone chứa điểm (Point-in-Polygon)}
Để biết điểm $p$ thuộc zone nào, hệ thống sử dụng \texttt{cv2.pointPolygonTest()} cho mỗi zone:

\begin{equation}
\texttt{pointPolygonTest}(Z_i, p) = \begin{cases}
    +d & \text{nếu } p \text{ nằm trong polygon } Z_i \text{ (khoảng cách đến biên)} \\
    0 & \text{nếu } p \text{ nằm trên biên} \\
    -d & \text{nếu } p \text{ nằm ngoài polygon}
\end{cases}
\end{equation}

Thuật toán dựa trên \textbf{ray casting}: vẽ tia từ $p$ ra vô cùng, đếm số lần cắt biên polygon --- lẻ = bên trong, chẵn = bên ngoài. Phức tạp $O(n)$ với $n$ cạnh polygon.

\subsection{Optical Flow --- Lucas-Kanade}
\label{sec:optical_flow}

Khi ball detector không phát hiện bóng, hệ thống sử dụng \textbf{Sparse Optical Flow} (phương pháp Lucas-Kanade \cite{lucas1981}) để ước tính vận tốc bóng từ chuyển động pixel.

\subsubsection{Giả thiết}
Phương pháp dựa trên 3 giả thiết:
\begin{enumerate}[noitemsep]
    \item \textbf{Brightness constancy}: Cường độ pixel không đổi giữa 2 frame: $I(x,y,t) = I(x+u, y+v, t+1)$
    \item \textbf{Small motion}: Chuyển vị $(u, v)$ nhỏ, cho phép khai triển Taylor bậc 1
    \item \textbf{Spatial coherence}: Các pixel lân cận có cùng vận tốc (cửa sổ $15 \times 15$)
\end{enumerate}

\subsubsection{Phương trình}
Khai triển Taylor và rút gọn:
\begin{equation}
I_x u + I_y v + I_t = 0
\end{equation}

trong đó $I_x, I_y$ là gradient spatial. Hệ thống dùng \textbf{đệ quy kim tự tháp} (pyramidal) với 3 tầng để xử lý chuyển động lớn:

\begin{equation}
\begin{bmatrix} u \\ v \end{bmatrix} = 
\left(\sum_{(x,y) \in W} \begin{bmatrix} I_x^2 & I_x I_y \\ I_x I_y & I_y^2 \end{bmatrix}\right)^{-1}
\sum_{(x,y) \in W} \begin{bmatrix} -I_x I_t \\ -I_y I_t \end{bmatrix}
\end{equation}

\subsubsection{Ứng dụng trong pipeline}
Tại vị trí bóng frame $t$, tìm vị trí tương ứng frame $t+1$ bằng optical flow:
\begin{equation}
(x_{t+1}, y_{t+1}) = (x_t + u, y_t + v)
\end{equation}

Vận tốc $(u, v)$ được blend vào trạng thái Kalman (xem Phần~\ref{sec:kalman_of}) với tỷ lệ $\alpha = 0.3$.

\subsection{Lọc vùng sân (Court Polygon Filter)}
\label{sec:court_polygon}

\subsubsection{Mục đích}
Loại bỏ người/vật ngoài sân (khán giả, trọng tài) trước khi phân tích. Chỉ giữ detections có \textbf{điểm chân} (foot point = bottom-center của bounding box) nằm trong polygon sân.

\subsubsection{Xây dựng polygon}
Từ 4 keypoints góc sân $(k_0, k_2, k_9, k_{11})$, tạo polygon \textbf{mở rộng 15\%} đều theo mọi hướng:

\begin{equation}
\bar{p} = \frac{1}{4} \sum_{i \in \{0,2,9,11\}} k_i \quad \text{(tâm polygon)}
\end{equation}

\begin{equation}
k'_i = k_i + \gamma \cdot (k_i - \bar{p}), \quad \gamma = 0.15
\end{equation}

Margin 15\% đảm bảo VĐV đứng sát biên vẫn được giữ lại, bù sai số keypoint detection.

\subsubsection{Foot point}
Điểm chân (contact point với mặt sân) ước tính từ bounding box:
\begin{equation}
\text{foot} = \left(\frac{x_1 + x_2}{2}, \; y_2 \right)
\end{equation}

trong đó $(x_1, y_1, x_2, y_2)$ là bbox coordinates (top-left, bottom-right). Dùng $y_2$ (đáy bbox) vì đó là vị trí chân trên mặt sân.

\subsection{Vẽ minimap (Court View Rendering)}
\label{sec:rendering}

Hệ thống render 3 bản đồ sân 2D riêng biệt, mỗi bản kích thước $150 \times 330$ pixels:

\subsubsection{Scale mapping}
Ánh xạ tọa độ sân 2D $(0\text{--}400, 0\text{--}880)$ sang pixel minimap:
\begin{equation}
px = \text{pad} + \text{round}(x_{2D} \cdot s_x), \quad py = \text{pad} + \text{round}(y_{2D} \cdot s_y)
\end{equation}

với $s_x = 150/400 = 0.375$, $s_y = 330/880 = 0.375$, pad $= 8$.

\subsubsection{Ball trail rendering}
Trail vẽ bằng \textbf{fade-in effect} --- các đoạn cũ mờ hơn, đoạn mới đậm hơn:
\begin{equation}
\alpha_k = \frac{k+1}{N_{trail}}, \quad \text{thickness}_k = \max(1, \lfloor 2\alpha_k \rfloor)
\end{equation}

\begin{equation}
\text{color}_k = \left(0, \; 140 + 115\alpha_k, \; 200 + 55\alpha_k\right) \quad \text{(xanh lá → vàng)}
\end{equation}

\subsubsection{Team color coding}
Vận động viên được mã hóa màu theo đội:
\begin{itemize}[noitemsep]
    \item Team A (gần camera): $(255, 150, 50)$ --- cam
    \item Team B (xa camera): $(50, 150, 255)$ --- xanh dương
    \item Bounce IN: $(0, 255, 0)$ --- xanh lá
    \item Bounce OUT: $(0, 0, 255)$ --- đỏ
\end{itemize}

% ============================================================
\section{Triển khai}
\label{sec:implementation}

\subsection{Môi trường}
Hệ thống được triển khai trên Google Colab với GPU NVIDIA A100/T4, sử dụng Python 3.12, PyTorch 2.x, Ultralytics (YOLO), OpenCV, FilterPy (Kalman), NumPy, và Pandas.

\textbf{Lý do sử dụng Colab}: Các mô hình YOLO11m yêu cầu GPU CUDA với VRAM tối thiểu 8GB cho inference và 16GB+ cho training. Máy local thông thường không đáp ứng yêu cầu này. Colab cung cấp GPU A100 (40GB) phù hợp cho cả training và inference.

\subsection{Trực quan hóa: Triple Minimap}
Hệ thống hiển thị \textbf{3 bản đồ sân 2D} xếp dọc:
\begin{enumerate}[noitemsep]
    \item \textbf{TRAIL}: Đường di chuyển bóng + vị trí vận động viên theo team
    \item \textbf{ALL BOUNCES}: Tất cả điểm bounce tích lũy (IN: xanh, OUT: đỏ)
    \item \textbf{LATEST}: Chỉ điểm bounce gần nhất (marker lớn, auto-clear)
\end{enumerate}

\subsection{Quá trình phát triển (7 phiên bản)}

\begin{table}[H]
\centering
\caption{Lịch sử phát triển pipeline qua 7 phiên bản}
\label{tab:versions}
\begin{tabularx}{\textwidth}{c X X}
\toprule
\textbf{Ver.} & \textbf{Tính năng mới} & \textbf{Vấn đề giải quyết} \\
\midrule
v1 & Court detection, ball tracking, homography, minimap & Framework ban đầu \\
v2 & Thêm ball bounding box, player detection (YOLOv8n), vẽ court lines & Cần hiển thị vị trí bóng và người chơi trên frame, dựng nền tảng phát hiện đa đối tượng \\
v3 & Kalman Filter cho ball trajectory, outlier rejection, proximity selection & Vị trí bóng nhảy lung tung (jitter) do model detect lệch vài pixel mỗi frame, và nhận nhầm vật thể tĩnh là bóng (false positive) \\
v4 & Half precision (FP16), giảm image size 640 & Tràn bộ nhớ GPU (Out of Memory) khi chạy trên A100 40GB \\
v5 & Bounce detection, Optical Flow fusion, In/Out classification & Cần xác định điểm bóng chạm sân và phân loại trong/ngoài sân \\
v6 & Zone-based projection: chia sân 6 zones, mỗi zone có local homography riêng & Global homography (1 ma trận cho toàn sân) gây sai lệch chiếu lớn ở vùng xa keypoints, đặc biệt góc sân \\
v7 & Court polygon filter loại người ngoài sân, Net-Y team assignment, cascade ball detection 4 tầng, triple minimap, bounce/hit filter, zigzag keypoint fix & 5 vấn đề: (1) phát hiện nhầm khán giả/trọng tài, (2) gán nhầm team khi detect 3 người, (3) tỷ lệ detect bóng thấp, (4) nhầm hit thành bounce, (5) minimap gương ngược do thứ tự keypoints \\
\bottomrule
\end{tabularx}
\end{table}

% ============================================================
\section{Kết quả thực nghiệm}
\label{sec:experiments}

\subsection{Player Detection}
Kết quả validation sau 80 epochs training:

\begin{table}[H]
\centering
\caption{Kết quả validation --- Player Detection Model (YOLO11m)}
\label{tab:results_player}
\begin{tabular}{lcccc}
\toprule
\textbf{Class} & \textbf{Precision} & \textbf{Recall} & \textbf{mAP@50} & \textbf{mAP@50-95} \\
\midrule
Person     & 0.968 & 0.980 & \textbf{0.982} & 0.796 \\
Pickleball & 0.846 & 0.761 & 0.797 & 0.349 \\
\midrule
\textbf{Overall} & 0.907 & 0.871 & 0.890 & 0.572 \\
\bottomrule
\end{tabular}
\end{table}

\begin{itemize}[noitemsep]
    \item Thời gian training: 1.25 giờ trên A100
    \item Inference speed: 1.7ms / image
    \item Person detection gần hoàn hảo (mAP@50 = 98.2\%)
    \item Pickleball detection khó hơn do kích thước nhỏ (mAP@50 = 79.7\%)
\end{itemize}

\subsection{Cascade Ball Detection}
So với phương pháp đơn mô hình (chỉ sử dụng tầng 1), cascade giúp tăng detection rate:

\begin{table}[H]
\centering
\caption{So sánh detection rate: đơn mô hình vs cascade}
\label{tab:cascade_compare}
\begin{tabular}{lcc}
\toprule
\textbf{Phương pháp} & \textbf{Detection Rate} & \textbf{False Positive Rate} \\
\midrule
YOLO ball model only & Thấp & Thấp \\
Cascade 4-stage & Cao hơn đáng kể & Tăng nhẹ (do classical CV) \\
\bottomrule
\end{tabular}
\end{table}

\textit{Ghi chú: Detection rate cụ thể phụ thuộc vào video input và được báo cáo qua cell diagnostics trong notebook.}

% ============================================================
\section{Khó khăn và thảo luận}
\label{sec:challenges}

\subsection{Giới hạn của Homography 2D}
Homography chỉ chính xác cho điểm trên mặt phẳng sân. Khi bóng bay cao, vị trí pixel camera không tương ứng vị trí trên sân --- gây sai lệch projection. Giải pháp: chỉ đánh giá IN/OUT tại thời điểm bounce.

\subsection{Phân biệt Bounce và Hit}
Cả bounce và hit đều gây đổi hướng bóng. Bộ lọc dựa trên vị trí tương đối bóng-player giảm false bounce, nhưng vẫn tồn tại edge case: bóng đập sát chân VĐV khi đang bị che khuất.

\subsection{Court Detection ở góc quay bất thường}
Model keypoint luôn output 12 keypoints, kể cả khi camera cận cảnh (chỉ thấy VĐV, không thấy sân). Hàm \texttt{is\_reliable()} lọc bằng confidence và diện tích, nhưng không thể xử lý tất cả trường hợp.

\subsection{Thiếu Player Re-Identification}
Hệ thống chưa có tracking ID liên tục giữa các frame. Team assignment dựa trên vị trí Y mỗi frame nên labels có thể thay đổi khi players di chuyển qua net.

% ============================================================
\section{Kết luận và hướng phát triển}
\label{sec:conclusion}

\subsection{Kết luận}
Bài báo đã trình bày hệ thống phân tích trận đấu pickleball end-to-end, phát triển qua 7 phiên bản với nhiều cải tiến kỹ thuật. Các đóng góp chính bao gồm:
\begin{enumerate}[noitemsep]
    \item Phát hiện zigzag pattern trong keypoint ordering và giải pháp mapping
    \item Cascade ball detection 4 tầng kết hợp deep learning và computer vision truyền thống
    \item Zone-based local homography cho projection chính xác hơn
    \item Bộ lọc bounce/hit dựa trên vị trí tương đối bóng-vận động viên
    \item Net-Y team assignment tránh lỗi ép chia đều
\end{enumerate}

\subsection{Hướng phát triển}
\begin{enumerate}[noitemsep]
    \item \textbf{Player Re-ID}: Tích hợp DeepSORT hoặc ByteTrack để tracking ID liên tục
    \item \textbf{3D Ball Tracking}: Sử dụng stereo camera hoặc depth estimation để biết chiều cao bóng
    \item \textbf{Action Recognition}: Phân loại loại đánh (serve, volley, dink, drive)
    \item \textbf{Real-time inference}: Tối ưu pipeline cho xử lý thời gian thực ($\geq 30$ FPS)
    \item \textbf{Multi-camera fusion}: Kết hợp nhiều góc camera để tăng coverage
\end{enumerate}

% ============================================================
\begin{thebibliography}{99}

\bibitem{sfia2023} Sports \& Fitness Industry Association, ``2023 Pickleball Participation Report,'' 2023.

\bibitem{hawkeye2024} Hawk-Eye Innovations, ``Ball Tracking Technology in Tennis,'' \url{https://www.hawkeyeinnovations.com}, 2024.

\bibitem{yolo2016} J. Redmon, S. Divvala, R. Girshick, and A. Farhadi, ``You Only Look Once: Unified, Real-Time Object Detection,'' in \textit{Proc. CVPR}, 2016.

\bibitem{ultralytics2024} Ultralytics, ``YOLO11 Documentation,'' \url{https://docs.ultralytics.com}, 2024.

\bibitem{faulkner2017} J. Faulkner, R. Dick, ``TennisCourt.net: A Deep Learning Approach to Tennis Court Detection and Tracking,'' 2017.

\bibitem{maji2022} D. Maji et al., ``YOLO-Pose: Enhancing YOLO for Multi Person Pose Estimation,'' in \textit{Proc. CVPRW}, 2022.

\bibitem{tracknet2019} Y.-C. Huang et al., ``TrackNet: A Deep Learning Network for Tracking High-speed and Tiny Objects in Sports Applications,'' in \textit{Proc. AVSS}, 2019.

\bibitem{kalman1960} R.E. Kalman, ``A New Approach to Linear Filtering and Prediction Problems,'' \textit{Trans. ASME--J. Basic Eng.}, vol.~82, pp.~35--45, 1960.

\bibitem{lucas1981} B. Lucas and T. Kanade, ``An Iterative Image Registration Technique with an Application to Stereo Vision,'' in \textit{Proc. DARPA IUW}, 1981.

\bibitem{hartley2003} R. Hartley and A. Zisserman, \textit{Multiple View Geometry in Computer Vision}, 2nd ed., Cambridge Univ. Press, 2003.

\bibitem{suzuki1985} S. Suzuki and K. Abe, ``Topological Structural Analysis of Digitized Binary Images by Border Following,'' \textit{Computer Vision, Graphics, and Image Processing}, vol.~30, no.~1, pp.~32--46, 1985.

\end{thebibliography}

\end{document}
